\page{09 / PASSAGE REVIEW}{Read overlap as a set of meanings}
\begin{takeaway}
\textbf{Review coverage.} The two Terra reviewers were assigned all 267 consolidated rows, all eight interviews, the corresponding original human records and all 985 AI claims. The lead read the calibration record and thematic narratives, inspected the returned judgments, and reopened contested evidence. The combined ledger records each row's interpretation, reasons and claim references.
\end{takeaway}

\subsection*{Three recurring relationships}
\begin{tabularx}{\linewidth}{@{}p{3.2cm}YY@{}}
\toprule
\textbf{Relationship} & \textbf{Example} & \textbf{What the comparison establishes}\tabularnewline\midrule
Same meaning, different names & T4 clan membership: cultural connection/preservation/traditions; AI organised Scottish meetings, dance and bagpipes. & Different labels can identify a common cultural-maintenance pattern. The book's divergence designation should not override that reading.\tabularnewline\addlinespace
Same passage, incomplete meaning & T6 the trunk and bundle containing all worldly goods; AI tracks the mother's management of children and luggage. & Practical caregiving is present, but material reduction to portable possessions needs a separate reading. A passage match does not establish completeness.\tabularnewline\addlinespace
Same meaning, different passage anchor & T3 a feverish infant removed from his mother; AI records the removal, uncertainty, three-day wait and reunion in several claims. & An unlinked human row can be substantively represented elsewhere in the same interview's AI analysis.\tabularnewline\bottomrule
\end{tabularx}

\subsection*{What the missed material adds}
Some residual omissions are consequential: Kaiz's derivative citizenship account makes legal dependence visible; severe journey illness adds embodied cost; housing without running water adds ordinary settlement deprivation; retrospective statements about childhood work can turn a list of chores into a life-course interpretation. These have different analytical value and should not each count as one identical ``miss.''

\subsection*{What the additional AI material adds}
Much of the unaligned AI material expands context rather than introducing a wholly new theme. The most useful contributions connect details: Larsen's furniture apprenticeship supplies blueprint skills for contracting; Thome's boarders and later rental flats reveal a changing household economy; Rizzuto's workplace exits connect wages and asserted freedom to an eventual restriction on mobility. Other claims merely subdivide a known episode. No overall novelty or precision rate is estimated for the 570 AI-only claims.

\begin{caution}
\textbf{The review required correction too.} Lead checks found that initial subagent judgments had mistaken missing lexical links for absent interpretations in several cases. Those judgments were reopened against the full claim inventory. This is why the report leads with inspectable cases and qualified comparative findings, rather than converting the ledger into a benchmark accuracy percentage.
\end{caution}
\source{Crosswalk rows 68, 100, 176; T1 derivative citizenship and journey illness; T8 housing; AI local craft, lodging and workplace themes. The ledger is a descriptive audit of model judgments, not independent human agreement.}

\page{10 / CLOSE-READING CASES}{Genuine omissions beside substantive recovery}
\case{T1 / Embodied transit costs remain underdeveloped}
{C1's ``Physical difficulties'' and C2/C3's physical-toll/difficult-journey readings are supported by severe sickness; a longer excerpt also includes a train-door hand injury. The meeting explicitly agrees on the shared hardship meaning.}
{The T1 claim set does not develop that sickness/injury experience. Claims about violence, separation and institutional vulnerability do not substitute for bodily suffering in transit. This is a supported omission. Rows 11 and 27 overlap, however, so they are not two independent missed events.}
{T1 source illness account; crosswalk rows 11 and 27; calibration lines 9--98; original human excerpts.}

\case{T1 / Legal dependency is a distinct human contribution}
{C2 codes ``Dependent legal status'' where Kaiz describes citizenship obtained through her father or husband. This makes a legal and gendered relationship visible.}
{No corresponding T1 AI claim develops derivative citizenship. Its household-authority theme does not capture this institutional relationship. This finding concerns what the narrator reports; the historical legal rule itself was not independently verified.}
{T1 opening citizenship account; crosswalk row 3; C2 initial code and explanation.}

\case{T2 / The AI does capture lasting distress}
{The narrator relates being under a train seat, dust and fear of suffocation to subsequent intolerance of enclosed spaces. Human trauma/lasting-impact codes are warranted as readings of her account.}
{The AI's ``Children's journey as freedom and ordeal'' explicitly preserves the enduring consequence, S297. This challenges the calibration's prediction that AI would necessarily miss such implicit trauma. It does not validate a clinical diagnosis or establish causation beyond the narrator's attribution.}
{T2 train journey; crosswalk row 46; AI local journey claim 7, S297; calibration at 1:17:11.}

\case{T3 / A lexical nonmatch is not an analytical absence}
{The mother spends three terrible days after officials take her feverish nursing infant. Human meanings include separation, health barriers and fear.}
{Aperture represents removal without explanation, fear of return, prolonged uncertainty, recovery and the mother's relief. Claims under family separation and the child's threshold jointly preserve the human meanings despite the row's missing lexical link.}
{Crosswalk row 68; AI T3 family-separation claims 9--11, S150--S153; child's-threshold claims 3--9; emotional-weight claims 2--3.}

\page{11 / INTERPRETIVE WARRANT}{Where the AI's explanation needs revision}
\case{Batta was 24 on arrival, not a child}
{T7's source header states birth in 1897 and arrival in 1921, ``AGE 24.'' The interview confirms the birth year. Her father left when she was three, but the eventual Ellis Island encounter occurred in adulthood.}
{The AI groups 14 T7 arrival/detention claims under ``Ellis Island as a child's threshold'' and treats meeting her unfamiliar father within a childhood-memory narrative. The individual events can be sound while the thematic developmental frame is wrong. Revise the scope to include adult arrival or explicitly separate Batta as a contrasting case.}
{T7 source lines 8--21 and 37--39; AI T7 child's-threshold block; childhood-memory claim 7, S311; family-separation claim 1 explicitly acknowledges adulthood.}

\case{Friedman's volunteering: bereavement, reciprocity and conjecture}
{C2's ``Volunteering as healing'' is supported in context: after her husband's death she felt lost, volunteered more, and says the work did a great deal for her. C1 also sees giving back to migrants.}
{Aperture captures translation, retained Russian, later service and personal benefit. Its theme title goes further: absence of arrival support \emph{prompting} volunteerism. The interviewer juxtaposes her service with the help her family lacked, but she does not directly identify that absence as the cause of her volunteering. Bereavement is the more explicit trigger. Retain the reparative interpretation as a hypothesis, not an established mechanism.}
{T5 source lines 218--236; crosswalk row 158; AI local volunteerism claims 2, 7 and 12, S522/S533/S551; theme definition.}

\case{A household economy is more than a list of tasks}
{Thome's human codes connect long hours, childhood chores, family responsibility and hardship. Some retrospective loss-of-childhood meaning remains less explicit in the AI.}
{The AI adds a concrete account: sixteen boarders, food rather than wages for the mother's labour, reliance on boarding when the father cannot work, and a later shift to rental income. This has analytical value. Yet causal shorthand in summaries and claims needs checking; a detailed ledger does not guarantee that every explanatory connection is supported.}
{T6 source boarding/work accounts; crosswalk rows 168, 173, 180, 185, 207--208; AI lodging claims 3, 7--8, 11--14; children's-household-labour claims 3--10.}

\begin{takeaway}
\textbf{The main risk is not merely hallucinated facts.} It is the movement from a supported event to a too-broad organising claim: adult experience becomes childhood experience; later service becomes a causal repair of earlier neglect; multiple local trajectories become a universal story. Human interpretations require the same scrutiny.
\end{takeaway}
